%% kapa: exam=1 round=23 year=2012 period=2 form=A subject=appraisal_law subject_ko=부동산관계법규 seq=27 number=67
%% source: Q-Net 감정평가사 기출문제 (cst00309) artlSeq=5052580 file=raw/1008/5052580/2012년 제23회 감정평가사 제1차시험 기출문제(2교시 A형).pdf
%% conversion: kordoc 4.24.51 -> tools/gen1.py (auto); review=proofread-sonnet; flags=none
\begin{question}{67}
국유재산법령상 일반재산에 관한 설명으로 옳은 것은?
\begin{choices}
  \item 일반재산을 개발하는 경우에 공공의 편익성이나 지역발전의 기여도 등은 고려해야 하지만 재정수입의 증대는 고려요소가 아니다.
  \item 용도를 지정하여 매각하는 일반재산은 그 재산의 매각일로부터 10년 이상 지정된 용도로 활용하여야 한다.
  \item 일반재산은 부동산신탁을 취급하는 신탁업자에게 신탁하여 개발하게 해서는 안 된다.
  \item 증권이 아닌 일반재산을 공공기관에 처분하는 경우에 예정가격은 두 개의 감정평가법인의 평가액을 산술평균한 금액으로 결정한다.
  \item 중앙관서의 장이 국가 외의 자가 소유하는 토지에 있는 국가소유의 건물을 그 토지 소유자에게 양여하는 경우에는 총괄청과 협의 없이 양여할 수 있다.
\end{choices}
\end{question}
\answer{2}
